\documentclass[letterpaper,11pt]{article}

\usepackage[T1]{fontenc}
\usepackage[utf8]{inputenc}
\usepackage[bitstream-charter]{mathdesign}
\usepackage[top=0.7in,bottom=0.7in,left=0.9in,right=0.9in]{geometry}
\usepackage[hidelinks]{hyperref}
\usepackage{parskip}
\linespread{1.04}
\setlength{\parskip}{7pt}
\pagestyle{empty}
\raggedright

\begin{document}

%----------HEADER----------
\begin{center}
    {\fontsize{18}{20}\selectfont\bfseries DEVA ANAND}\\ \vspace{3pt}
    {\small
    \href{mailto:devaanand@umass.edu}{devaanand@umass.edu}
    \enspace$\vert$\enspace
    (413) 315-1715
    \enspace$\vert$\enspace
    \href{https://devaanand.com}{devaanand.com}
    \enspace$\vert$\enspace
    \href{https://github.com/Deva-1903}{github.com/Deva-1903}}
\end{center}
\vspace{6pt}

\noindent July 15, 2026

\vspace{2pt}

\noindent PathAI\\
Boston, MA

\vspace{2pt}

\noindent \textbf{Re: Machine Learning Intern/Co-op (Fall 2026)}

\vspace{4pt}

\noindent Dear PathAI Team,

The first research I ever published was a deep-learning model that classified disease from medical images, so PathAI's work, turning pathology images into more accurate diagnoses for patients with cancer, is remarkably close to the problem that first pulled me into machine learning. I am an MS Computer Science student at UMass Amherst, and I would love to spend this fall contributing to models that genuinely change patient outcomes.

That first-author IEEE paper applied deep transfer learning to neuroimaging for medical image classification, and it taught me early that in a clinical setting the science has to be rigorous or it does not matter. I have kept that habit since. In a recent interpretability study I measured a safety behavior in an LLM using difference-in-means probing, bootstrap 95\% confidence intervals, and McNemar's tests rather than eyeballing results, and as the sole engineer on a UMass research lab's ML platform I built a model registry and per-run reproducibility exports so any analysis a researcher runs can be reproduced exactly. Careful, reproducible analysis is not overhead to me, it is the point.

I also enjoy the full lifecycle of a model, not just training it. In one project I took a base model end to end through supervised fine-tuning, DPO, and GRPO, raised held-out accuracy from 54\% to 72\%, then served it behind a vLLM and FastAPI gateway with caching and micro-batching and watched throughput and latency on Grafana. That mix of model development plus the MLOps and pipeline work around it is exactly what your intern description covers, and I am comfortable moving between the two. My two years as a backend engineer at Wysa, a clinical mental-health company, and my current work on a clinical decision support tool also mean I am used to building software where patient impact is real.

I will be upfront that my computer-vision depth is more medical-imaging and coursework than digital pathology, but that is precisely the gap I want to close on a team doing rigorous science at this scale, and I am the kind of person who is glad to take on any task, big or small, to help the team ship. Thank you for considering my application, and for building something that matters.

\vspace{6pt}

\noindent Sincerely,\\
Deva Anand

\end{document}
